% !TeX root = ../navier-stokes-manuscript.tex
% ============================================================
% 2.1 Vortex Stretching
% ============================================================

The momentum equation has four pieces:
\[
  \underbrace{\partial_t u}_{\text{local acceleration}}
  +
  \underbrace{(u\cdot\nabla)u}_{\text{nonlinear transport}}
  =
  \underbrace{-\nabla p}_{\text{pressure}}
  +
  \underbrace{\nu\Delta u}_{\text{viscous diffusion}}.
\]
The central regularity problem lies in the competition between nonlinear transport and viscosity.
The hope is simple: viscosity smooths the fluid faster than nonlinear transport can sharpen it.
If that remains true at every scale, the velocity stays smooth.

\subsection{Vortex Stretching}\label{sub:ThePerfectlyStretchingVortex}

In a narrow material vortex tube, the vorticity points along the tube's axis.
When the surrounding strain draws one segment out in that direction, the segment lengthens.
Its volume cannot change, so it narrows at the same time.
If \(e\) is the axial direction and \(L(t)\) is the segment's length, the symmetric strain
\[
  S:=\frac12\bigl(\nabla u+(\nabla u)^{\mathsf T}\bigr),
\]
records the extension through
\[
  Se=\sigma(t)e,
  \qquad
  \sigma(t)>0,
  \qquad
  \frac{\dot L(t)}{L(t)}
  =
  e\cdot Se
  =
  \sigma(t).
\]
The same extension fixes the transverse contraction.
Writing the constant material volume as \(\mathcal V=A(t)L(t)\), with effective radius \(r_{\mathrm{eff}}(t):=\sqrt{A(t)/\pi}\), incompressibility gives
\[
  \nabla\cdot u=0,
  \qquad
  \frac{d}{dt}(A L)=0,
  \qquad
  \frac{\dot A}{A}=-\sigma(t),
  \qquad
  \frac{\dot r_{\mathrm{eff}}}{r_{\mathrm{eff}}}
  =
  -\frac{\sigma(t)}{2}.
\]
The tube is being driven toward a smaller spatial scale as it is pulled out along its axis.

Its rotation changes in the same motion.
For \(\omega:=\nabla\times u\),
\[
  \frac{D\omega}{Dt}
  =
  S\omega+\nu\Delta\omega,
  \qquad
  \frac{D}{Dt}
  :=
  \partial_t+u\cdot\nabla,
\]
When \(\omega=|\omega|e\), the stretching direction and the vorticity direction coincide, and
\[
  S\omega
  =
  \sigma(t)\omega,
  \qquad
  \frac12\frac{D|\omega|^2}{Dt}
  =
  \sigma(t)|\omega|^2
  +
  \nu\,\omega\cdot\Delta\omega.
\]
Where this complete right-hand side remains positive, the segment's rotation intensifies while its radius falls.
That growth already sits inside the velocity gradient.
If \(|u|\leq U_0\) on the material segment \(\Omega_{\mathrm{seg}}(t)\), then
\[
  E_{\mathrm{seg}}(t)
  :=
  \frac12\int_{\Omega_{\mathrm{seg}}(t)}|u(x,t)|^2\,dx
  \leq
  \frac12U_0^2\mathcal V
  <
  \infty,
  \qquad
  \left|\frac12\bigl(\nabla u-(\nabla u)^{\mathsf T}\bigr)\right|_{\mathrm F}
  =
  \frac{|\omega|}{\sqrt2}
  \leq
  |\nabla u|_{\mathrm F}.
\]
The vortex can therefore remain finite in kinetic energy while becoming longer, thinner, and more sharply rotating.
The gradients rise.

\divider{\NavierStokes{} Scaling}

The falling radius has made scale part of the problem.
At a fixed time \(t_0\), let \(\ell:=r_{\mathrm{eff}}(t_0)\), and for \(\lambda>1\) define
\[
  u_\lambda(x,t)
  :=
  \lambda u(\lambda x,\lambda^2t),
  \qquad
  p_\lambda(x,t)
  :=
  \lambda^2p(\lambda x,\lambda^2t).
\]
At time \(\lambda^{-2}t_0\), the corresponding vortex in \(u_\lambda\) has width \(\lambda^{-1}\ell\).
This compares solutions at different scales; it does not identify \(u_\lambda\) with a later state of the original material tube.
The four terms transform together:
\[
\begin{aligned}
  \partial_tu_\lambda(x,t)
  &=
  \lambda^3(\partial_tu)(\lambda x,\lambda^2t),
  \\
  (u_\lambda\cdot\nabla)u_\lambda(x,t)
  &=
  \lambda^3\bigl((u\cdot\nabla)u\bigr)(\lambda x,\lambda^2t),
  \\
  \nabla p_\lambda(x,t)
  &=
  \lambda^3(\nabla p)(\lambda x,\lambda^2t),
  \\
  \nu\Delta u_\lambda(x,t)
  &=
  \lambda^3\nu(\Delta u)(\lambda x,\lambda^2t).
\end{aligned}
\]
The common \(\lambda^3\) is the answer: the finer width has strengthened viscosity and nonlinear sharpening together.

\divider{Critical Height}

The equation has kept its balance, but the size assigned to the velocity depends on where it is measured.
With \(\Lambda:=(-\Delta)^{1/2}\), a change of variables gives
\[
  \norm{u_\lambda(t)}_{\dot H^s}^2
  =
  \lambda^{2s-1}
  \norm{u(\lambda^2t)}_{\dot H^s}^2.
\]
Below \(s=\tfrac12\) the norm falls under concentration; above it the norm rises.
At \(s=\tfrac12\), the change of scale contributes nothing at all.
Set
\[
  H(t)
  :=
  \frac12\norm{u(t)}_{\dot H^{1/2}}^2
  =
  \frac12\norm{\Lambda^{1/2}u(t)}_2^2.
\]
The levels on either side move in opposite directions:
\[
  \norm{u_\lambda(t)}_2^2
  =
  \lambda^{-1}\norm{u(\lambda^2t)}_2^2,
  \qquad
  H_\lambda(t)=H(\lambda^2t),
  \qquad
  \norm{\nabla u_\lambda(t)}_2^2
  =
  \lambda\norm{\nabla u(\lambda^2t)}_2^2.
\]
Since concentration leaves \(H\) unchanged even as the neighboring levels separate, any subsequent motion of \(H\) must come from the equation itself, not from the shrinking ruler.

\divider{Nonlinear Production versus Viscous Dissipation}

Applying \(\Lambda^{1/2}\) to the momentum equation and pairing with \(\Lambda^{1/2}u\) isolates that change.
The signed nonlinear contribution is
\[
  \mathcal P_{1/2}(t)
  :=
  -\operatorname{Re}
  \pair
    {\Lambda^{1/2}u(t)}
    {\Lambda^{1/2}\bigl((u(t)\cdot\nabla)u(t)\bigr)}_{L^2}.
\]
Incompressibility removes the pressure pairing, and the Laplacian contributes \(-\nu\norm{\Lambda^{3/2}u}_2^2\), leaving
\[
  H'(t)
  +
  \nu\norm{\Lambda^{3/2}u(t)}_2^2
  =
  \mathcal P_{1/2}(t).
\]
Nonlinear production raises the critical height; viscous dissipation lowers it.
Under the same rescaling,
\[
  \mathcal P_{1/2}[u_\lambda](t)
  =
  \lambda^2\mathcal P_{1/2}[u](\lambda^2t),
  \qquad
  \nu\norm{\Lambda^{3/2}u_\lambda(t)}_2^2
  =
  \lambda^2\nu\norm{\Lambda^{3/2}u(\lambda^2t)}_2^2.
\]
The common \(\lambda^2\) leaves the balance untouched and moves the finer comparison onto the contracted coordinate \(t\mapsto\lambda^{-2}t\).

\divider{The Heat Clock}

Viscosity carries exactly such a clock.
For the linear heat equation \(v_t=\nu\Delta v\),
\[
  \widehat v(\xi,t+\delta t)
  =
  e^{-\nu|\xi|^2\delta t}\widehat v(\xi,t).
\]
A structure of width \(r\) is concentrated around frequencies \(|\xi|\simeq r^{-1}\), and the multiplier changes by order one once \(\nu|\xi|^2\delta t\simeq1\).
Its viscous comparison time is therefore
\[
  \tau_\nu(r)
  :=
  \frac{r^2}{\nu}.
\]
Shrinking the width by \(\lambda^{-1}\) gives
\[
  \tau_\nu(\lambda^{-1}r)
  =
  \lambda^{-2}\tau_\nu(r),
\]
the same contraction already carried by the full \NavierStokes{} scaling.
This is the time viscosity needs to alter structure at width \(r\); it is not a prescribed lifetime for the original vortex.
Even so, continued narrowing would place each finer scale beside a quadratically shorter clock.

\divider{The Zeno Cascade}

Repeated halving makes those clocks a geometric sequence.
For
\[
  r_n:=2^{-n}r_0,
  \qquad
  \tau_n:=\frac{r_n^2}{\nu},
\]
one has
\[
  \tau_n
  =
  \frac{r_0^2}{\nu}4^{-n},
  \qquad
  \sum_{n=0}^{\infty}\tau_n
  =
  \frac{4r_0^2}{3\nu}
  <
  \infty.
\]
Infinitely many comparison clocks fit inside a finite total time.
The sum does not make them one fluid history.
A candidate singularity would require the same solution to pass through widths
\[
  r_0>r_1>r_2>\cdots\longrightarrow0
\]
at times
\[
  t_0<t_1<t_2<\cdots\uparrow T_*<\infty,
  \qquad
  t_{n+1}-t_n\asymp\tau_n,
\]
with comparison constants independent of \(n\).
Its remaining time would then satisfy
\[
  T_*-t_n
  \asymp
  \sum_{k=n}^{\infty}\tau_k
  =
  \frac{4r_n^2}{3\nu}.
\]
The width tends to zero inside intervals of the same order as its viscous clock.
For one solution to make that passage, its response must keep changing on those collapsing intervals.

Assume that this same solution loses continuation at \(T_*\) through blow-up in one fixed continuation-controlling Banach space \(B\).
If some finite temporal row \(\partial_t^N u\), \(N\geq1\), remained uniformly bounded in \(B\) after a time \(t_0<T_*\), then
\[
  \partial_t^{N-1}u(t)
  =
  \partial_t^{N-1}u(t_0)
  +
  \int_{t_0}^{t}\partial_t^N u(s)\,ds
\]
would remain bounded because the interval is finite.
The same identity applied successively to the lower rows would bound \(u\), contradicting the assumed blow-up.
No finite temporal row can therefore retain uniform control in \(B\) along a genuine singular history.

The momentum equation supplies the first row, and each further time differentiation supplies the next.
The resulting family
\[
  u,
  \qquad
  \partial_tu,
  \qquad
  \partial_t^2u,
  \qquad
  \partial_t^3u,
  \qquad
  \ldots
\]
is the temporal Tower.
